\documentclass[10pt]{article}
\usepackage{authblk, amssymb, amsmath, amsthm, latexsym, array, enumerate}


%\oddsidemargin=0cm
%\topmargin=-1,5cm
%\textwidth=16cm
%\textheight=23cm
%\hfuzz=1.111pt

\usepackage[english,russian]{babel}
\usepackage[utf8]{inputenc} % sic!
\usepackage[a5paper,margin=2cm,innermargin=2cm]{geometry}

\makeatletter
\renewenvironment{thebibliography}[1]
{\subsubsection*{\refname}%
	\@mkboth{\MakeUppercase\refname}{\MakeUppercase\refname}%
	\list{\@biblabel{\@arabic\c@enumiv}}%
	{\settowidth\labelwidth{\@biblabel{#1}}%
		\leftmargin\labelwidth
		\advance\leftmargin\labelsep
		\@openbib@code
		\usecounter{enumiv}%
		\let\p@enumiv\@empty
		\renewcommand\theenumiv{\@arabic\c@enumiv}}%
	\sloppy
	\clubpenalty4000
	\@clubpenalty \clubpenalty
	\widowpenalty4000%
	\sfcode`\.\@m}
{\def\@noitemerr
	{\@latex@warning{Empty `thebibliography' environment}}%
	\endlist}
\makeatother


%%%%%%%%%%%%%%%%%%%%%%%%%%%%%%%

{\theoremstyle {definition} \newtheorem {defi} {Definition[section] }
	{\theoremstyle {plain}  \newtheorem {theorem} [defi] {Theorem}}
	{\theoremstyle {plain}  \newtheorem {corollary} [defi]{Corollary}}
	{\theoremstyle {plain}  \newtheorem {lemma} [defi]{Lemma}}
	{\theoremstyle {plain}  \newtheorem {proposition} [defi]{Proposition}}
	{\theoremstyle {plain} \newtheorem {remark}[defi] {Remark}}
	{\theoremstyle {definition} \newtheorem {example}[defi] {Example}}
	{\theoremstyle {definition} \newtheorem {problem}[defi] {Problem}}
	{\theoremstyle {definition} \newtheorem {definition}[defi] {Definition}}

%%%%%%%%%%%%%%%%%%%%%%%%%%%%%%%
%%%%%%%%%%%%%%%%%%%%%%%%%%%%%%%
%%%%%%%%%%%%%%%%%%%%%%%%%%%%%%%

\title{\Large 
	Вопросы адаптации шага по времени \\ в модели электрического пробоя \\ типа диффузной границы
%\footnote{Информация о финансировании.}
}

\author[1]{\underline{Андрей Пономарев}}

\affil[1]{ИПМ им. М. В. Келдыша РАН}

\date{}

\renewcommand*{\Affilfont}{\small\it}
\renewcommand*{\Authfont}{\bf\normalsize}

\makeatletter
\renewcommand{\Authand}{, }      % вместо "and" ставим запятую
\renewcommand{\Authands}{, }     % для списка авторов
\renewcommand{\Authsep}{, }      % разделитель между авторами
\makeatother

\begin{document}

\maketitle


Предмет исследования автора -- модель развития канала электрического пробоя в твердом диэлектрике, относящаяся к типу \linebreak диффузной границы, предложенная в статье \cite{ponopitike}. Одна из сложностей при работе с моделью состоит в существенно различных временных масштабах процессов: как правило, электрический пробой развивается стремительно, но ему предшествует долгий период медленных изменений в системе. Доклад посвящен возможному решению проблемы -- адаптации расчетного шага по времени в модели.

Автором исследовано три различных метода адаптации шага по~времени, их работа проверена численно для задачи в упрощенной, одномерной, постановке. Подробный анализ методов приведен в~статье~\cite{ponoadaptation}.

При выборе диапазона допустимых шагов по времени следует принимать во внимание несколько ограничений. Используемая явная разностная схема требует выполнения условий устойчивости~\cite{ponostability}. Также шаг по времени должен быть кратно меньше характерного времени релаксации заряда в среде.

Наконец, самый удачный из исследованных методов адаптации внедрен в разрабатываемую научным коллективом программу \linebreak для~моделирования задачи в полной постановке. Проведены расчеты программы в~двух модельных постановках: развитие канала пробоя из~проводящей иголки и из случайных малых повреждений образца.

Отметим, что в результате внедрения адаптации реальное время расчета снизилось кратно.

{
\begin{thebibliography}{9}

\bibitem{ponopitike}
K.~Ch.~Pitike, W.~Hong, {\it Phase-field model for dielectric breakdown in~solids}, Journal of Applied Physics {\bf 115(4)}, 044101 (2014).

\bibitem{ponoadaptation}
А.~С.~Пономарев [и др.], {\it Адаптация шага по времени в~модели типа <<диффузной границы>> на основе уравнения Аллена--Кана}, Математическое моделирование {\bf 38(4)}, 3--22 (2026).

\bibitem{ponostability}
А.~С.~Пономарев, Е.~В.~Зипунова, Е.~Б.~Савенков, {\it Устойчивость стационарных решений в модели развития канала электрического пробоя типа <<диффузной границы>>}, Препринты ИПМ им.~М.~В.~Келдыша {\bf 72}, 1--23 (2024).

\end{thebibliography}
}

\end{document}
